% QMB_n07_chsh_experiment.tex — Chapter 7: ΔCHSH and Experiment
% Sources: Doc. 190, 230, 147
\chapter{CHSH Deviation and Experimental Reach}
\label{ch:chsh_exp}

\section{Two Different Observables}

The correction register (Doc.~190, Precision~11 and 12, May 2026)
distinguishes precisely between two different quantities that earlier
documents conflated:

\textbf{Elementary deviation.} Per individual measurement,
phase costs of the $\theta_4$ projection cycle arise. The natural
candidate is:
\begin{equation}
  \Delta\mathrm{CHSH}_{\mathrm{elem}} = \frac{\xi}{2\pi}
  \approx 2{,}1\times10^{-5}.
\end{equation}
This identification is a working hypothesis; a formal derivation
from the $T^4$ projection geometry is still missing.\status{S}

\textbf{Cumulative deviation.} Over $N$ measurement runs, the
deviation grows according to the formula from Doc.~022:
\begin{equation}
  \Delta\mathrm{CHSH}(N) = 2\sqrt2 \Bigl[1 - \exp\!\Bigl(-\xi\,\frac{\ln N}{D_f}\Bigr)\Bigr]
  \approx \xi\,\frac{\ln N}{D_f}\cdot2\sqrt2.
\end{equation}
For $N=73$ this yields $\approx5\times10^{-4}$. This quantity is
experimentally accessible once error correction thresholds are
surpassed.

\section{Experimental Situation 2026}

IBM Heron~r2 (May 2026, Doc.~147) provides:
\begin{itemize}
  \item Bell violation confirmed: $S=2{,}74$ ($96{,}9\,\%$ of the
        Tsirelson bound).\status{K}
  \item FFGFT circuits operational on current hardware.\status{K}
  \item Device-to-device variation: $1{,}4\,\%$ — approximately
        $1400\times$ larger than the sought $\xi$ effect.\status{K}
\end{itemize}

The $\xi$ effect requires sub-ppm-stable error-corrected qubits.
Not resolvable with NISQ hardware.\status{K}

\section{Open Derivation}

The identification $\Delta\mathrm{CHSH}_{\mathrm{elem}} = \xi/(2\pi)$
is plausible, but not derived from the $T^4$ projection geometry.
The ratio between cumulative and elementary deviation for
$N$ runs is:
\begin{equation}
  \frac{\Delta\mathrm{CHSH}(N)}{\Delta\mathrm{CHSH}_{\mathrm{elem}}}
  \approx N\cdot\frac{2\pi}{D_f},
\end{equation}
for $N=73$ a factor $\approx153$, consistent with the observed
factor $\sim50$--$100$.\status{S}

The formal derivation of $\xi/(2\pi)$ from the carrier geometry
remains an open bridge.

\quelle{Doc.~022, 147, 190 (Prec.~11, 12), 230}